% TOP_PROBLEM: {"id": 3405, "title": "P vs. BPP", "category_id": 15, "category": {"id": 15, "name": "computer_science", "display_name": "Computer Science", "description": "Computational complexity, algorithms, and theoretical CS.", "slug": "computer-science", "order_index": 15, "created_at": "2026-07-31T15:26:25.671Z"}, "status": "open"}
% FIRST_POSTED: 2026-09-25
% ENTRY_KIND: statement_only
% !TeX program = lualatex
% Definition and sources only; this file is not an LLM solution attempt.
\documentclass[11pt]{article}
\usepackage[margin=25mm]{geometry}
\usepackage{fontspec}
\setmainfont{Latin Modern Roman}
\usepackage{amsmath,amssymb,mathtools}
\usepackage{unicode-math}
\setmathfont{Latin Modern Math}
\usepackage{xurl,hyperref}
\hypersetup{colorlinks=true,urlcolor=blue}
\setcounter{secnumdepth}{0}
\setlength{\parindent}{0pt}
\setlength{\parskip}{5pt}
\setlength{\emergencystretch}{3em}
\begin{document}
\section*{\texorpdfstring{P vs. BPP}{P vs. BPP}}\label{3405}
\subsection{Definitions and mathematical statement}
A language $L$ is in $\mathsf{BPP}$ if a randomized polynomial-time algorithm $A$ satisfies $\Pr[A(x)=1_L(x)]\ge2/3$ for every input $x$. Membership in $\mathsf P$ requires a deterministic polynomial-time decision algorithm with no errors. Determine whether
\[
 \forall L\in\mathsf{BPP},\quad L\in\mathsf P.
\]
The simulation must work in the ordinary uniform model, not merely through nonuniform advice or for most inputs.
\par\smallskip\textit{Formulation and concept references:} \href{https://people.seas.harvard.edu/~salil/pseudorandomness/pseudorandomness-Aug12.pdf}{[S1]}.

\subsection{Short English statement}
Can randomness always be removed from an efficient bounded-error decision algorithm without losing polynomial-time efficiency?
\subsection{Sources}
\begin{itemize}
\item[S1]S. Vadhan, 'Pseudorandomness', Foundations and Trends in Theoretical Computer Science 7(1-3):1-336, 2012.
\url{https://people.seas.harvard.edu/~salil/pseudorandomness/pseudorandomness-Aug12.pdf}.
\textit{Formulation source.}
\item[S2]CS 6810: Theory of Computing, Lecture 15, March 17, 2026.
\url{https://www.cs.cornell.edu/courses/cs6810/2026sp/lec15.pdf}.
\textit{Further reference; not a proof endorsement.}
\item[S3]Pseudorandomness Beating the Hybrid Argument for Insensitive Algorithms.
\url{https://eccc.weizmann.ac.il/report/2026/082/}.
\textit{Further reference; not a proof endorsement.}
\end{itemize}
\end{document}
